\documentclass[10pt,a4paper]{amsart}
\usepackage[margin=19mm]{geometry}
\usepackage{amsmath,amssymb,mathtools}
\usepackage[scr=rsfs,cal=euler]{mathalfa}
\usepackage{xfrac}
\usepackage[unicode,hidelinks,bookmarks=false]{hyperref}
\newcommand{\ie}{i.e.\ }
\newcommand{\cf}{cf.\ }
\newcommand{\resp}{respectively\ }
\newcommand{\Subsection}{\subsection}
\providecommand{\nobreakdash}{\nobreak\hbox{-}}
\setlength{\emergencystretch}{2em}
\multlinegap0pt
\usepackage{bm}
\usepackage{bbm}
\usepackage{dsfont}
\usepackage{xspace}
\usepackage{cancel}
\usepackage{mathtools}
\usepackage{amsthm}
\makeatletter
\@ifundefined{theorem}{\newtheorem{theorem}{Theorem}}{}
\makeatother
\makeatletter
\expandafter\ifx\csname sketch\endcsname\relax
\newenvironment{sketch}{\begin{proof}[Sketch of proof.]}{\end{proof}}
\fi
\makeatother
\title[Log-Concavity and Log-Convexity of Restricted Infinite Produ]{Log-Concavity and Log-Convexity of Restricted Infinite Products}
\author{Krystian Gajdzica, Bernhard Heim, Markus Neuhauser}
\date{Source: arXiv:2509.11246v1 (CC BY 4.0)}
\begin{document}
\maketitle

\noindent\emph{Source and licence.} Log-Concavity and Log-Convexity of Restricted Infinite Products. Version arXiv:2509.11246v1 (CC BY 4.0), distributed under the Creative Commons Attribution 4.0 International licence, as recorded in the arXiv licence metadata for this exact version and shown on the abstract page https://arxiv.org/abs/2509.11246v1. Full rights notice: licenses/arxiv-2509.11246-CC-BY-4.0.txt.\par\smallskip

\noindent\textit{Editorial scope.} Counted unit: the Theorem whose source label is \texttt{thm: involved} in the article (source0.tex lines 1204--1223, source label \texttt{thm: involved}), delivered together with the article's own complete original proof of it (source0.tex lines 1225--1295, one \texttt{proof} environment of 71 lines). No mathematics is written, repaired or re-derived: statement and proof are the article's own text line for line. Required context retained uncounted: def: A\_E(n) (lines 1197--1199). Every definition, display and label of those lines is retained unchanged. Omitted from this package and not redistributed: 1--1196, 1200--1203, 1224--1224, 1296--1670. Nothing inside a retained excerpt is omitted or altered: every retained body is delivered exactly as the article's lines are written, so no omission receipt of any kind accompanies this package. Citations: the retained excerpts carry no citation command at all, so no bibliography entry of the article is transcribed and no reference is redistributed. Reference adaptation: none was needed and none was made; every reference inside the delivered unit resolves inside it, so no unresolved reference or citation is produced. Printed slips kept literal: no printed word, symbol, hypothesis or number of the counted statement or of its proof was altered. Layout and class: the article's own class is \texttt{amsart} with packages including amssymb, amsmath, amsthm, comment, pdfpages, graphicx, placeins, wrapfig, float, lineno, youngtab, setspace, multirow; the wrapper is the lane's pinned \texttt{amsart} editorial wrapper at 10pt on A4, which restates the theorem-like environments the retained text needs and adds the article's own macro definitions that the retained text needs (none), with extra wrapper packages (bm, bbm, dsfont, xspace, cancel, mathtools). The delivered numbering is the unit's own. Assets: no retained span contains a figure environment, a graphics-inclusion command, a PDF-page inclusion, a table or a TikZ picture; no image file is copied into this package, the package declares no asset dependency and no image file is redistributed with it. Licence: version arXiv:2509.11246v1 of this article is distributed under the Creative Commons Attribution 4.0 International licence, as recorded in the arXiv licence metadata for this exact version and shown on the abstract page https://arxiv.org/abs/2509.11246v1. A full rights notice is delivered at licenses/arxiv-2509.11246-CC-BY-4.0.txt. Characters: every retained excerpt is pure ASCII, so the delivered engine, XeLaTeX with its default Latin Modern text font, reports no missing character.
\begin{align}\label{def: A_E(n)}
A_E(n)\prod_{i=1}^k\frac{g_{E,\ell }(m_i)}{m_i}\leq p_{E,\ell }(n) \leq A_E(n)\prod_{i=1}^k\frac{g_{E,\ell }(m_i)}{m_i}+p_{E,1}(n)\Tilde{M}_E(n)^{\Psi(\ell)},
\end{align}

\begin{theorem}\label{thm: involved}
Let $E$, $\Phi$, $\Psi$, and
$\{f_{\ell}\}_{\ell}$ be given. Suppose further that the inequality
$M_E(n)>M_E(n-1)$ holds for every $n\geq n_0$ and some positive integer
$n_0$. Let us also assume that $Q_{E}\left( n\right) =1$ for some
$n>n_0$, and  the maximal products $M_E(n), M_E(n-1)$,
and $M_E(n+1)$ are attained by the unique $S$-partitions $(x_1,\ldots,x_a)$,
$(y_1,\ldots,y_b)$, and $(z_1,\ldots,z_c)$ (respectively) such that
\begin{align*}
\frac{g_{\ell}(x_1)^2\cdots g_{\ell}(x_a)^2}{x_1^2\cdots x_a^2}\sim
\frac{g_{\ell}(y_1)\cdots g_{\ell}(y_b)g_{\ell}(z_1)\cdots g_{\ell}(z_c)}{y_1\cdots y_bz_1\cdots z_c}
\end{align*}
for large $\ell $.
Then $p_{E,\ell }(n)$ is strictly
\begin{enumerate}
\item  log-concave at $n$ for all sufficiently large values of $\ell$ if $\frac{A_E(n)^2}{A_E(n-1)A_E(n+1)}>1$,
\item  log-convex at $n$ for all sufficiently large values of $\ell$ if $\frac{A_E(n)^2}{A_E(n-1)A_E(n+1)}<1$,
\end{enumerate}
where the values $A_E(n), A_E(n-1)$, and $A_E(n+1)$ are defined as in  (\ref{def: A_E(n)}).
\end{theorem}

\begin{proof}
Since both cases are very similar, let us only deal with the criterion for the log-concavity property. The assumptions from the statement together with the inequalities (\ref{def: A_E(n)}) assert that
\begin{equation*}
p_{E,\ell }(n)^2\geq A_E(n)^2\left(\prod_{i=1}^a\frac{g_{E,\ell }(x_i)}{x_i}\right)^2
\end{equation*}
and
\begin{align*}
p_{E,\ell }(n-1) \, p_{E,\ell }(n+1)\leq &\left(A_E(n-1)\prod_{i=1}^b\frac{g_{E,\ell }(y_i)}{y_i}+p_{E,1}(n-1)\Tilde{M}_E(n-1)^{\Psi(\ell)}\right)\\
&{}\cdot
\left(A_E(n+1)\prod_{i=1}^c\frac{g_{E,\ell }(z_i)}{z_i}+p_{E,1}(n+1)\Tilde{M}_E(n+1)^{\Psi(\ell)}\right)\\
<& A_E(n-1)A_E(n+1)\prod_{i=1}^b\frac{g_{E,\ell }(y_i)}{y_i}\prod_{i=1}^c\frac{g_{E,\ell }(z_i)}{z_i}\\
&{}+
A_E(n-1)p_{E,1}(n+1)M_E(n-1)^{\Psi(\ell)+1}\Tilde{M}_E(n+1)^{\Psi(\ell)}\\
&{}+
A_E(n+1)p_{E,1}(n-1)M_E(n+1)^{\Psi(\ell)+1}\Tilde{M}_E(n-1)^{\Psi(\ell)}\\
&{}+
p_{E,1}(n-1)p_{E,1}(n+1)\left(\Tilde{M}_E(n-1)\Tilde{M}_E(n+1)\right)^{\Psi(\ell)}.
\end{align*}
Next, we want to compare all of the above four summands with the estimations for $p_{\ell,S}(n)$.
In the first case, we obtain that
\begin{align*}
\frac{A_E(n-1)A_E(n+1)\prod_{i=1}^b\frac{g_{E,\ell }(y_i)}{y_i}\prod_{i=1}^c\frac{g_{E,\ell }(z_i)}{z_i}}{A_E(n)^2\left(\prod_{i=1}^a\frac{g_{E,\ell }(x_i)}{x_i}\right)^2}
\sim \frac{A_E(n-1)A_E(n+1)}{A_E(n)^2}<1
\end{align*}
for sufficiently large values of
$\ell $.
The second one, on the other hand, requires a little bit more computations:
\begin{align*}
&\frac{A_E(n-1)p_{E,1}(n+1)M_E(n-1)^{\Psi(\ell)+1}\Tilde{M}_E(n+1)^{\Psi(\ell)}}{A_E(n)^2\left(\prod_{i=1}^a\frac{g_{E,\ell }(x_i)}{x_i}\right)^2}\\
\leq
&2\frac{A_E(n-1)p_{E,1}(n+1)M_E(n-1)^{\Psi(\ell)+1}\Tilde{M}_E(n+1)^{\Psi(\ell)}}{A_E(n)^2\prod_{i=1}^b\frac{g_{E,\ell }(y_i)}{y_i}\prod_{i=1}^c\frac{g_{E,\ell }(z_i)}{z_i}}\\
\leq& 2\frac{A_E(n-1)p_{E,1}(n+1)M_E(n-1)^{\Psi(\ell)+1}\Tilde{M}_E(n+1)^{\Psi(\ell)}}{A_E(n)^2\left(M_E(n-1)M_E(n+1)\right)^{\Phi(\ell)}}\\
\leq
&2\frac{A_E(n-1)p_{E,1}(n+1)M_E(n-1)^{\Psi(\ell)+1}\Tilde{M}_E(n+1)^{\Psi(\ell)}}{A_E(n)^2\left(M_E(n-1)M_E(n+1)\right)^{-
B}\left(M_E(n-1)M_E(n+1)\right)^{\Psi(\ell)}}\\
=&C_1(n)\left(\frac{\Tilde{M}_E(n+1)}{M_E(n+1)}\right)^{\Psi(\ell)}\rightarrow 0
\end{align*}
for sufficiently large $\ell $.
The third case is similar to the second one and presents as follows:
\begin{align*}
&\frac{A_E(n+1)p_{E,1}(n-1)M_E(n+1)^{\Psi(\ell)+1}\Tilde{M}_E(n-1)^{\Psi(\ell)}}{A_E(n)^2\left(\prod_{i=1}^a\frac{g_{E,\ell }(x_i)}{x_i}\right)^2}\\
\leq
&2\frac{A_E(n+1)p_{E,1}(n-1)M_E(n+1)^{\Psi(\ell)+1}\Tilde{M}_E(n-1)^{\Psi(\ell)}}{A_E(n)^2\prod_{i=1}^b\frac{g_{E,\ell }(y_i)}{y_i}\prod_{i=1}^c\frac{g_{E,\ell }(z_i)}{z_i}}\\
\leq& 2\frac{A_E(n+1)p_{E,1}(n-1)M_E(n+1)^{\Psi(\ell)+1}\Tilde{M}_E(n-1)^{\Psi(\ell)}}{A_E(n)^2\left(M_E(n-1)M_E(n+1)\right)^{\Phi(\ell)}}\\
\leq
&2\frac{A_E(n+1)p_{E,1}(n-1)M_E(n+1)^{\Psi(\ell)+1}\Tilde{M}_E(n-1)^{\Psi(\ell)}}{A_E(n)^2\left(M_E(n-1)M_E(n+1)\right)^{-
B}\left(M_E(n-1)M_E(n+1)\right)^{\Psi(\ell)}}\\
=&C_2(n)\left(\frac{\Tilde{M}_E(n-1)}{M_E(n-1)}\right)^{\Psi(\ell)}\rightarrow 0
\end{align*}
for sufficiently large $\ell $.
The last instance can be estimated in the following way
\begin{align*}
&\frac{p_{E,1}(n-1)p_{E,1}(n+1)\left(\Tilde{M}_E(n-1)\Tilde{M}_E(n+1)\right)^{\Psi(\ell)}}{A_E(n)^2\left(\prod_{i=1}^a\frac{g_{E,\ell }(x_i)}{x_i}\right)^2}\\
\leq
&2\frac{p_{E,1}(n-1)p_{E,1}(n+1)\left(\Tilde{M}_E(n-1)\Tilde{M}_E(n+1)\right)^{\Psi(\ell)}}{A_E(n)^2\prod_{i=1}^b\frac{g_{E,\ell }(y_i)}{y_i}\prod_{i=1}^c\frac{g_{E,\ell }(z_i)}{z_i}}\\
\leq& 2\frac{p_{E,1}(n-1)p_{E,1}(n+1)\left(\Tilde{M}_E(n-1)\Tilde{M}_E(n+1)\right)^{\Psi(\ell)}}{A_E(n)^2\left(M_E(n-1)M_E(n+1)\right)^{\Phi(\ell)}}\\
\leq
&2\frac{p_{E,1}(n-1)p_{E,1}(n+1)\left(\Tilde{M}_E(n-1)\Tilde{M}_E(n+1)\right)^{\Psi(\ell)}}{A_E(n)^2\left(M_E(n-1)M_E(n+1)\right)^{-
B}\left(M_E(n-1)M_E(n+1)\right)^{\Psi(\ell)}}\\
=&C_3(n)\left(\frac{\Tilde{M}_E(n-1)\Tilde{M}_E(n+1)}{M_E(n-1)M_E(n+1)}\right)^{\Psi(\ell)}\rightarrow 0
\end{align*}
for sufficiently large values of
$\ell $.
All of the above calculations lead us finally to the conclusion that
\begin{align*}
\frac{p_{E,\ell }(n-1)p_{E,\ell }(n+1)}{p_{E,\ell }(n)^2}\rightarrow \frac{A_E(n-1)A_E(n+1)}{A_E(n)^2}<1,
\end{align*}
for sufficiently large values of
$\ell $
and this finishes the proof.
\end{proof}

\end{document}
